\begin{figure}[ht]
\centering
\begin{tikzpicture}[>=Stealth, font=\small]
  \draw[->] (-4.2,0) -- (4.4,0) node[right] {$\Re(s)$};
  \draw[->] (0,-2.4) -- (0,2.6) node[above] {$\Im(s)$};
  \draw[dashed, thick] (1,-2.2) -- (1,2.4);
  \node[above right] at (1,2.2) {$\Re(s)=1$};
  \draw[very thick, blue!70!black, ->] (2.8,1.4) .. controls (2.0,1.7) and (1.4,1.7) .. (1.02,1.4);
  \draw[very thick, blue!70!black, ->] (-0.8,1.4) .. controls (-0.2,1.7) and (0.5,1.7) .. (0.98,1.4);
  \draw[very thick, blue!70!black, ->] (2.8,-1.4) .. controls (2.0,-1.7) and (1.4,-1.7) .. (1.02,-1.4);
  \draw[very thick, blue!70!black, ->] (-0.8,-1.4) .. controls (-0.2,-1.7) and (0.5,-1.7) .. (0.98,-1.4);
  \fill[red!70!black] (1,0) circle (2pt);
  \node[below right] at (1,0) {中心点 $s=1$};
  \node[fill=white, inner sep=2pt] at (2.45,1.85) {$s\mapsto 2-s$};
\end{tikzpicture}
\caption{权 $2$ 模形式完成的 $L$-函数关于直线 $\Re(s)=1$ 满足函数方程对称。中心点 $s=1$ 正是 Birch--Swinnerton-Dyer 猜想讨论零点阶数的位置。}
\label{fig:ch08-functional-equation}
\end{figure}
